\chapter{统计学习理论与模型选择}

\begin{quote}\small
\textit{本章摘要。}本章为模型比较提供统计学习视角。先说明假设空间和容量如何影响泛化界，再讨论结构风险、交叉验证、校准、不平衡数据和置信区间。最后把理论约束转成数据泄漏检查，使模型选择从“哪个分数最高”转向“哪个结论在何种不确定性下成立”。
\end{quote}

上一章说明训练误差、抽样误差和分布偏移不是一回事，但现实实验还多了一层选择：研究者会试很多窗口、特征和超参数，再报告其中最好的结果。这时需要控制的不再是某个预先固定函数的误差，而是整个选择过程可能带来的乐观偏差。本章以这个区别为出发点，把统计学习的容量约束与实际验证规则连接起来。

试过许多窗口长度以后再挑最高分，相当于又用数据作了一次选择。若仍用同一批数据报告最终成绩，就可能把偶然的好运当成稳定的提升。理论中的复杂度代价说明这种问题为什么会发生，交叉验证帮助安排有限样本的用途，校准和成本分析则帮助把分数变成告警规则。每一种方法都有适用条件，不能替代对新设备的独立检验。

阅读本章前，应能区分训练平均损失与总体风险，并理解一次数据划分为何不能证明所有工况下的表现。第一次阅读可先掌握有限候选集合的联合界和嵌套验证；无限类复杂度推导用于回答连续参数模型如何控制选择偏差，不是实施交叉验证的前置操作。

\section{假设空间与泛化约束}

% auxiliary-resource-guide
本章末的“辅助学习材料”列出与核心方法对应的讲解和实践入口，可在阅读相关推导后，按所列小节对照学习。

\begin{figure}[htbp]
\centering
\begin{tikzpicture}[x=1cm,y=0.75cm]
 \draw[->] (0,0)--(7.2,0) node[right] {模型复杂度};
 \draw[->] (0,0)--(0,4.0) node[above] {误差};
 \draw[blue!70!black, thick] plot[smooth] coordinates {(0.5,3.5)(1.5,2.6)(2.5,1.8)(3.5,1.25)(4.5,0.95)(5.5,0.8)(6.8,0.72)};
 \draw[red!70!black, thick] plot[smooth] coordinates {(0.5,1.0)(1.5,1.25)(2.5,1.7)(3.5,2.2)(4.5,2.9)(5.5,3.35)(6.8,3.7)};
 \draw[black!65, thick, dashed] plot[smooth] coordinates {(0.5,3.7)(1.5,2.9)(2.5,2.2)(3.5,1.75)(4.5,1.55)(5.5,1.65)(6.8,2.05)};
 \node[blue!70!black, font=\scriptsize] at (5.4,0.55) {训练误差};
 \node[red!70!black, font=\scriptsize] at (5.7,3.7) {复杂度代价（示意）};
 \node[font=\scriptsize] at (3.9,1.95) {真实风险};
 \draw[dashed] (4.35,0)--(4.35,1.5);
 \node[font=\scriptsize, anchor=north] at (4.35,-0.08) {折中点};
\end{tikzpicture}
\caption{模型复杂度同时影响拟合误差与估计方差；模型选择的目标是控制真实风险，而不是单独追求训练误差。}
\label{fig:statistical-complexity-tradeoff}
\end{figure}
图\ref{fig:statistical-complexity-tradeoff}是定性示意，不是风险分解的数值相加，也不保证真实风险总呈U形或只有一个最低点。它提醒读者，选模应控制未见样本风险，而非仅追求更低训练误差；理论复杂度上界也不等于实测的估计方差。
\paragraph{假设空间。}
学习算法从假设空间$\mathcal{H}$中选择函数。有限训练集只能约束$\mathcal{H}$在观测点上的行为，因此假设空间越丰富，插值训练数据越容易，但泛化分析也越困难。模型容量不能简单等同于参数数量，还与参数化方式、数据几何和优化偏置有关。

\paragraph{读概率界前先确定谁在随机变化。}
假设空间$\mathcal H$是一组允许的预测函数，而不是数据表。以标量温度$x$为输入，三个预先固定的阈值分类器$f_c(x)=\mathbb 1\{x>c\}$、$c\in\{60,70,80\}$就组成含三个函数的空间。若阈值可以取任意实数，则空间无限，尽管每个函数只有一个参数。

以下$\Pr$指重新抽取训练集时的概率，$\epsilon>0$是允许的风险误差，$\delta$是允许理论保证失败的概率。$\sup$表示上确界，有限集合时就是最大值；$|\mathcal H|$是函数个数，$\log$是自然对数。固定函数的界像保证一个预先选定候选；同时界像保证整张候选名单，随后再看数据选人也仍在名单内。它不是说每个样本的预测误差都小于$\epsilon$。

\paragraph{统一收敛的直观证明。}
对固定函数$f$，若损失$\ell(f(X),Y)\in[0,1]$，Hoeffding不等式给出
\begin{equation}
 \Pr\left(|R(f)-\hat R(f)|>\epsilon\right)\leq2\exp(-2n\epsilon^2).
\end{equation}
若$\mathcal{H}$是有限集合，对所有$f\in\mathcal{H}$使用联合界可得
\begin{equation}
 \Pr\left(\sup_{f\in\mathcal{H}}|R(f)-\hat R(f)|>\epsilon\right)
 \leq2|\mathcal{H}|\exp(-2n\epsilon^2).\label{eq:statistical-union-bound}
\end{equation}
这说明样本量需要随假设空间复杂度增加。对无限空间，可用 Vapnik--Chervonenkis (VC) 维、Rademacher复杂度或覆盖数替代$|\mathcal{H}|$。

\paragraph{从概率尾界到样本量。}
以上$n$个样本必须独立同分布，$\mathcal H$应预先固定且非空；平方损失无界时不能不加条件地使用式\eqref{eq:statistical-union-bound}。为看到指数项的来源，令$Z_i=\ell(f(X_i),Y_i)\in[0,1]$。有界变量的指数矩满足$\E\exp(t(\E Z_i-Z_i))\leq\exp(t^2/8)$；独立性使指数矩相乘。Markov不等式给出
\begin{equation}
 \Pr(R-\hat R>\epsilon)
 \leq\exp(-tn\epsilon+nt^2/8).
\end{equation}
取$t=4\epsilon$得到单侧$\exp(-2n\epsilon^2)$，对另一侧重复并相加即得双侧界。联合界不要求各函数的误差独立，因此适用于在同一数据上比较的多个模型。

令式\eqref{eq:statistical-union-bound}右端为$\delta\in(0,1)$，解得以至少$1-\delta$的概率，对所有$f\in\mathcal H$同时成立
\begin{equation}
 |R(f)-\hat R(f)|\leq b_n,
 \qquad b_n=\sqrt{\frac{\log(2|\mathcal H|/\delta)}{2n}}.\label{eq:statistical-uniform-radius}
\end{equation}
若希望误差不超过$\epsilon$，充分条件是$n\geq\log(2|\mathcal H|/\delta)/(2\epsilon^2)$。这是一条保守的充分保证，不是对具体任务所需数据量的精确预测。

设$\hat f$最小化经验风险，$f_{\mathcal H}^*$最小化类内真实风险，则在式\eqref{eq:statistical-uniform-radius}对所有函数同时成立的事件上
\begin{align}
 R(\hat f)&\leq\hat R(\hat f)+b_n
 \leq\hat R(f_{\mathcal H}^*)+b_n
 \leq R(f_{\mathcal H}^*)+2b_n.
\end{align}
第一步和第三步都需要同时界；只给预先固定$f$的界无法直接代入数据依赖的$\hat f$。若优化仅达到$\eta$精度，最后再加$\eta$。相对所有可测函数中最优风险的差还包含类外近似误差，故增大$n$不能弥补错误的函数类。

\paragraph{代入数字检验界的含义。}
若$|\mathcal H|=10$、$n=1000$、$\delta=0.05$，则$b_n=\sqrt{\log(400)/2000}\approx0.055$。经验错误率为$0.08$的所选函数，在同时界成立时真实错误率不超过约$0.135$；相对类内最优函数的保证则需$2b_n\approx0.110$，不能把这两个误差半径混用。若要求$b_n\leq0.05$，充分样本数为$\lceil\log(400)/(2\times0.05^2)\rceil=1199$。这些数值较保守，说明“理论能证明”与“实验已经达到”是两种证据。

无限类段落中的VC维可先用一维阈值理解：一个点可以任意标为零或一，但两个由小到大排列的点不能得到“一、零”的标记，因此该类的VC维为一。Rademacher随机符号则相当于独立抛硬币生成的正负记号；复杂度询问函数类能在多大程度上追随这种无规律波动。读者可先掌握这两个对象，再阅读它们的上界推导。

\paragraph{结构风险最小化。}
经验风险最小化只看训练误差，结构风险最小化在一系列嵌套假设空间中选择复杂度。实际的正则化、树深度限制、网络宽度和数据增强都可视为控制有效复杂度的方法。理论上更强的模型仍可能泛化良好，因为优化算法偏爱某些简单或平滑的解。

\paragraph{复杂度惩罚如何随候选类分配。}
对预先指定的每个有限类$\mathcal H_m$，选择正权重$\pi_m$。权重之和为一，并据此分配失败概率：
\begin{equation*}
 \sum_m\pi_m=1,\qquad \delta_m=\delta\pi_m.
\end{equation*}
再次使用联合界，得到所有类同时有效的偏差界
\begin{equation}
 b_m=\sqrt{\frac{\log(2|\mathcal H_m|/(\delta\pi_m))}{2n}}.
\end{equation}
若在所有$m,f\in\mathcal H_m$中最小化$\hat R(f)+b_m$，则所选模型$\tilde f$满足
\begin{equation}
 R(\tilde f)\leq\inf_m\inf_{f\in\mathcal H_m}\{R(f)+2b_m\}.
\end{equation}
推导是先用$R(\tilde f)\leq\hat R(\tilde f)+b_{\tilde m}$，再用最小化性质与每个候选比较，最后把候选的经验风险换成真实风险。复杂类要靠真实的拟合改善抵消较大的惩罚，而不是仅靠训练误差获胜。

\paragraph{无限类需要什么替代量。}
连续参数类通常有无限多个函数，不能把参数个数代入$\log|\mathcal H|$。二分类的增长函数$\Pi_{\mathcal H}(n)$计数在$n$个点上最多能实现多少种二元标记；若VC维为$d_{\rm VC}$，则$n\geq d_{\rm VC}\geq1$时有$\Pi_{\mathcal H}(n)\leq(en/d_{\rm VC})^{d_{\rm VC}}$。它衡量可区分的标记，而非参数存储量。

另一个视角是随机符号拟合能力。令损失类$\mathcal G=\{(x,y)\mapsto\ell(f(x),y):f\in\mathcal H\}$，独立符号$\epsilon_i$等概率为$\pm1$，定义
\begin{equation}
 \mathfrak R_n(\mathcal G)=\E_{D,\epsilon}\sup_{g\in\mathcal G}
 \frac1n\sum_i\epsilon_i g(X_i,Y_i).
\end{equation}
这里另记$Z_i=(X_i,Y_i)$为一条样本，而非前面Hoeffding推导中的标量损失。引入独立同分布的复制样本$D'=\{Z_i'\}_{i=1}^n$后，用$R(g)=\E_{D'}\hat R_{D'}(g)$及上确界的凸性，得到
\begin{align}
 \E_D\sup_g(R(g)-\hat R_D(g))
 &\leq\E_{D,D'}\sup_g\frac1n\sum_i(g(Z_i')-g(Z_i))\\
 &\leq2\mathfrak R_n(\mathcal G).
\end{align}
第二步利用样本对可交换性插入随机符号，再将两个上确界分开。这解释了“能拟合纯随机波动”为什么意味着更大的乐观风险。若损失在$[0,1]$内，替换一个样本使上确界至多改变$1/n$，有界差分集中进一步给出以至少$1-\delta$的概率，$\sup_g(R-\hat R)\leq2\mathfrak R_n(\mathcal G)+\sqrt{\log(1/\delta)/(2n)}$。这些结论仍依赖独立抽样和适当可测性，并非对重叠窗口的自动保证。

\section{模型选择、校准与不平衡}
模型选择偏差要求把调参与最终评价分离\cite{cawley2010selection}。深度网络校准研究则表明，分类正确率与概率可信度需要分别检验，温度缩放是一种后处理基线\cite{guo2017calibration}。
\paragraph{模型选择与交叉验证。}
若数据近似独立同分布，$K$折交叉验证将训练集分为$K$份，轮流使用一份验证、其余训练。时间序列必须使用滚动窗口或前向验证，避免未来信息泄漏。超参数搜索也应嵌套在训练数据内部；否则测试集会逐渐变成隐形训练集。

\paragraph{交叉验证估计的是哪个算法。}
设折索引为$I_1,\ldots,I_K$，超参数为$h$，$\hat f_h^{(-k)}$表示只用第$k$折以外的数据完成预处理和训练后的模型。加权交叉验证风险为
\begin{equation}
 \widehat R_{\rm CV}(h)=\frac1n\sum_{k=1}^K\sum_{i\in I_k}
 \ell(\hat f_h^{(-k)}(x_i),y_i).
\end{equation}
每个损失的标签没有参与产生该预测，但不同折的训练集重叠，因此这些损失不是独立试验。选$\hat h=\argmin_h\widehat R_{\rm CV}(h)$后，最小分数有选择偏差：对带噪估计$\hat r_h$，总有$\E\min_h\hat r_h\leq\min_h\E\hat r_h$。搜索越广，越容易挑中有利噪声。

嵌套验证在外层训练集内部完成全部内层选择，再对外层留出样本预测，评估的是“选择加训练”这一完整算法。最后可用所有开发数据重新选择并拟合，但外层分数并非这个最终参数向量的无偏测量。设备切分估计跨设备能力，前向切分估计未来能力；两者对应不同目标，不能由普通随机$K$折替代。

\paragraph{一轮嵌套选择到底留下了什么。}
例如有六台开发设备，外层先留下设备5、6；只在设备1至4内部比较两个窗口长度，每个内层训练折都重新拟合填补和标准化。选好长度后，用设备1至4重新训练，最后才预测设备5、6。外层换组后，内层可能选到不同长度，这正是完整选择算法的表现，而不是执行错误。开发结束后另有独立测试设备时，应冻结最终流程再测试。超参数是训练前需选择的设置（如窗口长度），模型参数则是给定设置后从训练数据估计的数值（如回归系数）。

\paragraph{校准与决策。}
分类器输出的分数不一定是概率。校准要求在所有预测为$p$的样本中，正例频率接近$p$。可靠性图和 Expected Calibration Error 可用于检查校准。若预测要驱动检修、放行或停机决策，还应定义效用函数$U(a,y)$，最终选择
\begin{equation}
 a^*(x)=\argmax_a\mathbb{E}_{p(y\mid x,\mathcal{D})}[U(a,y)].
\end{equation}
这说明“最准确的分类器”不一定产生“最优的行动”。

\paragraph{不平衡与置信区间。}
少数类任务应报告每类召回率、宏平均$F_1$和 Area Under the Precision--Recall Curve (PR-AUC)。重采样、类别权重和阈值调整针对的是不同环节，不能混为一谈。性能估计还需要不确定性，可使用自助法在测试样本上重采样得到指标分布。跨设备评估时，设备级重采样比样本级重采样更能反映真实部署波动。

\paragraph{先读懂混淆矩阵，再读少数类指标。}
把故障记为正类，TP、FP、FN分别为正确告警、错误告警和漏报的数量。召回率为$\mathrm{TP}/(\mathrm{TP}+\mathrm{FN})$，精确率为$\mathrm{TP}/(\mathrm{TP}+\mathrm{FP})$，$F_1=2\mathrm{TP}/(2\mathrm{TP}+\mathrm{FP}+\mathrm{FN})$。例如100台设备中有10台故障，告警12台且其中8台故障，则召回率为$0.8$、精确率为$2/3$、$F_1=16/22\approx0.727$。一直报正常也有$90\%$准确率，却漏掉全部故障。若没有真实正例，召回率分母为零；若没有正类预测，精确率分母为零，此时原比值无定义，应说明报告约定，不能默认为性能完美。宏平均是先分别计算各类指标再等权平均，而不是先合并计数。PR曲线扫描阈值形成精确率与召回率的对应关系；面积还应声明插值或平均精确率等计算约定。

\paragraph{类别权重为何改变概率含义。}
令真实条件正例概率为$p$，预测为$q$，正负类权重为$a,b>0$。固定输入处的加权交叉熵为$L(q)=-ap\log q-b(1-p)\log(1-q)$。求导并令其为零，得到
\begin{equation}
 -\frac{ap}{q}+\frac{b(1-p)}{1-q}=0,
 \qquad q^*=\frac{ap}{ap+b(1-p)}.
\end{equation}
所以加权训练的理想输出一般不再是原始分布的$p$。若模型恰好达到该最优且权重已知，可反解$p=bq/[a(1-q)+bq]$；有限样本和错设模型下还需独立校准。重采样同样可能改变类别先验；它不应与只在预测后调整阈值混为一谈。

例如$p=0.1,a=9,b=1$时，理想加权输出为$q^*=0.5$，并不表示真实故障率升到了50\%；代入反解仍得到$p=0.1$。

精确率也直接受基率影响。若正例比例为$\pi$、真正率为$u$、假正率为$v$，Bayes公式给出
\begin{equation}
 \operatorname{Precision}=\frac{\pi u}{\pi u+(1-\pi)v}.
\end{equation}
即使排序和条件错误率不变，部署正例更稀少时精确率也会下降。跨工况比较精确率--召回率（Precision--Recall，PR）曲线应同时报告基率。

\paragraph{区间与配对比较的抽样单位。}
若独立同分布测试单位上的损失$L_i$具有有限方差，令$\bar L=n^{-1}\sum_iL_i$、$s_L^2=(n-1)^{-1}\sum_i(L_i-\bar L)^2$（$n>1$）。$s_L$描述单个单位的损失波动，均值的标准误估计$s_L/\sqrt n$描述重新抽取测试集时平均损失的波动，两者不能混用。大样本近似区间为$\bar L\pm1.96s_L/\sqrt n$，其中1.96是标准正态分布的97.5\%分位数，左右尾各留2.5\%。例如$n=100,\bar L=4,s_L=2$时，标准误为0.2，区间为$[3.608,4.392]$。这只针对固定模型和既定测试分布，不包括训练随机性；稀有事件或小样本时正态近似可能不良。

比较模型A、B时应先在同一测试单位上计算$d_i=L_{A,i}-L_{B,i}$，再估计$\bar d$的区间，因为
\begin{equation}
 \Var(L_A-L_B)=\Var(L_A)+\Var(L_B)-2\operatorname{Cov}(L_A,L_B).
\end{equation}
忽略配对会丢失共同设备难度的信息。对$G$台近似独立设备，自助法每次有放回抽取$G$个设备编号，保留所抽设备的全部窗口并重算指标；重复$B$次后取经验分位数作为近似区间。每台等权与每窗口等权是不同的目标，必须在每次重采样中维持原定义。设备数很少、设备之间相关或许多重采样没有正例时，区间本身也不可靠；增加重采样次数不会增加独立设备信息。

\paragraph{置信区间的概率究竟修饰什么。}
设固定模型的总体平均损失为$\mu$，由随机测试数据$D$算出的区间为$[L(D),U(D)]$。95\%置信区间的目标是$P_D\{L(D)\leq\mu\leq U(D)\}\approx0.95$：重复采集测试集并构造区间，约95\%的区间覆盖固定的$\mu$。一次数据已经到手后，不能仅凭这一定义把它读成“$\mu$有95\%概率位于当前两端点之间”。后续贝叶斯章节的可信区间则在给定数据和先验后，对未知参数的后验概率作陈述。预测区间又针对一个未来观测，不是总体均值；三者的对象和随机性不同。

自助法不是制造新设备。例如原有设备编号为$(A,B,C)$，一次有放回抽样可能为$(C,C,A)$，必须把$C$的全部记录按两次抽中计入，而非去重；另一次可能为$(B,A,B)$。每次重算相同指标，所得波动近似原实验的抽样波动。只重采样测试设备、不重新训练时，估计的是固定模型的测试不确定性，不包含重训和选模波动。三个设备仅用于演示操作，不足以支撑可靠的95\%区间。

\section{从统计结论到实验设置}
\paragraph{数据泄漏清单。}
常见泄漏包括：用全数据计算标准化参数，把同一设备的相邻窗口分到训练和测试，利用故障发生后的变量预测故障发生前的状态，以及根据测试集反复选择模型。泄漏往往带来异常漂亮的结果，因此每个实验都应记录切分单位、时间边界、预处理拟合范围和标签生成时刻。

\paragraph{从固定函数到数据依赖的模型选择。}
Hoeffding 界直接控制固定函数的经验风险偏差，但实际算法会根据数据选择函数。联合界的作用是同时控制整个假设空间，使“先看数据再选模型”仍有可分析的上界。代价是模型空间越复杂，需要更多样本或更强的结构约束。

\paragraph{容量不是参数数量的同义词。}
参数数量相同，网络能表达的函数也未必相同。激活函数、连接方式和参数共享决定函数类，优化路径则影响训练最后选到其中哪些函数。例如卷积在不同位置使用同一组权重，注意力允许相距较远的输入直接交互。分析泛化时，因此还要看函数类、参数范数、分类间隔和算法稳定性，不能仅按参数个数给模型排队。

\paragraph{交叉验证的工程细节。}
预处理、特征选择和超参数搜索都必须在每个训练折内部拟合。若先对全体数据做标准化或选择相关特征，再执行交叉验证，验证折的信息已经泄漏。时间序列使用前向验证时，验证窗口必须位于训练窗口之后，必要时在二者之间设置间隔以避免标签窗口重叠。

\paragraph{不平衡、阈值与区间估计。}
ROC-AUC 衡量排序能力，PR-AUC 更关注少数正类下的精确率和召回率。阈值改变会改变混淆矩阵，却不改变排序本身。若最终动作依赖阈值，应报告阈值扫描、成本曲线和置信区间。自助法重采样时，设备级任务应按设备重采样；否则同一设备的多个窗口会让区间过窄。

\paragraph{校准的后处理。}
温度缩放在校准集上用一个正标量调整 logits（归一化前的类别得分），保持同一样本内部的类别次序，但不保证改善独立测试集的校准；Platt scaling 用逻辑函数映射分数，isotonic regression（保序回归）拟合非递减映射，二者对数据量和函数形状的要求不同。校准模型不能使用测试集反复调参，且分布变化后需要重新检查。对安全告警而言，概率校准和拒识机制常比少数百分点的准确率更重要。

\paragraph{温度缩放的目标与边界。}
先冻结基模型及其预处理，再在校准集上保存logits并只拟合温度。部署时沿用相同变换；若用校准标签继续更新基模型，这批数据便不再是该模型的独立校准集，需要重新安排划分。
对$C$类logits $z_i\in\R^C$，在独立校准集上令$T>0$，定义$q_{ic}(T)=\exp(z_{ic}/T)/\sum_j\exp(z_{ij}/T)$，最小化
\begin{equation}
 L(T)=\sum_i\left[-z_{i,y_i}/T+\log\sum_c\exp(z_{ic}/T)\right].
\end{equation}
取逆温度$\beta=1/T$，则
\begin{align}
 L'(\beta)&=\sum_i\left(\sum_cq_{ic}z_{ic}-z_{i,y_i}\right),\\
 L''(\beta)&=\sum_i\Var_{c\sim q_i}(z_{ic})\geq0.
\end{align}
求导时使用$\partial_\beta q_{ic}=q_{ic}(z_{ic}-\sum_jq_{ij}z_{ij})$，故二阶导数正是得分的加权方差。因此这是一个一维凸优化问题（约束$\beta>0$）；下确界也可能只在$\beta\to0$或$\beta\to\infty$时逼近，不能保证存在有限最优温度。正温度保持每个样本内部的类别次序和最大类别；多分类中不同样本的最大概率排序则未必保持，不能据此保证所有排序指标不变。它只能整体调整尖锐程度，无法修复所有类别和工况上的系统偏差。

二分类可靠性图把样本分入$B_m$，计算箱内预测均值$\bar q_m$和事件比例$\bar y_m$；常用误差为$\sum_m(|B_m|/N)|\bar q_m-\bar y_m|$。这是依赖分箱的经验统计量，不是证明条件校准的定理。箱数太少会掩盖局部偏差，太多则产生高方差，应同时报告样本量与可靠性图。

\paragraph{自测：样本数怎样改变保证。}
固定候选集合和失败概率，若独立样本数从$n$增至$4n$，统一界半径$b_n$如何变化？把每条旧样本复制四次是否等价？

\textit{参考解。}式\eqref{eq:statistical-uniform-radius}给出$b_{4n}=b_n/2$；经验风险最小化相对类内最优的$2b_n$项也减半。复制不产生新的独立观测，不能把复制后的行数代入这一保证。

\section{从统计量到决策}
\paragraph{效用与损失只差一个决策方向。}
效用$U(a,y)$越大越好，损失$C(a,y)$越小越好，取$U=-C$便得到同一动作。若不检修时故障损失为100、无故障损失为零；检修固定花费8且在这个简化例子中消除该次故障损失，则故障概率为$0.1$时不检修的期望损失为10，检修为8，故选检修。这里必须写明检修效果；若只检测而不消除风险，需要把后续动作和残余损失也列入成本，不能直接沿用阈值。

模型输出只是决策系统的一个中间量。若动作集合为$\mathcal A$，效用函数为$U(a,y)$，最优动作是期望效用最大的动作。不同部门可能面对不同成本，因此同一个概率模型可以配合不同阈值，而不必为每种成本重新训练。这个分离有助于审计模型能力和业务策略各自的变化。

统计学习中的复杂度界与实验估计可结合官方交叉验证文档和重采样教程复习，注意它们服务于不同的不确定性来源。
\section*{辅助学习材料}
\begingroup
\emergencystretch=3em
\begin{itemize}
\item \textbf{Cross-validation: evaluating estimator performance}；作者/平台：scikit-learn Developers；英文；官方文档。阅读 K-fold、group、stratified 与 nested selection 的说明，帮助落实本章的模型选择和泛化估计。\href{https://scikit-learn.org/stable/modules/cross\_validation.html}{在线阅读}。
\item \textbf{Prediction Intervals for Gradient Boosting Regression}；作者/平台：scikit-learn Developers；英文；官方示例。阅读分位数回归、区间覆盖率与调参部分，区分预测区间和参数置信区间；该示例不提供分布无关的覆盖保证。\href{https://scikit-learn.org/stable/auto\_examples/ensemble/plot\_gradient\_boosting\_quantile.html}{在线阅读}。
\item \textbf{【机器学习】什么是交叉验证？为什么以及如何进行交叉验证？}；知乎；中文解说。区分模型选择与最终泛化估计；标题与主题据必应索引，原页受限。\href{https://zhuanlan.zhihu.com/p/81738187251}{在线阅读}。
\item \textbf{交叉验证方法汇总【附代码】（索引标题节选）}；CSDN；中文博客。比较留一法、K折与分层划分，另行检查设备分组要求；标题与主题据必应索引，原页受限，完整标题未核实。\href{https://blog.csdn.net/WHYbeHERE/article/details/108192957}{在线阅读}。
\item \textbf{2022 - 再探宝可梦、数码宝贝分类器 — 浅谈机器学习原理}；李宏毅；bilibili；中文课程。选看第23集，结合本章假设空间与泛化条件讨论分类。2026年10月4日官方公开元数据显示整套课程累计播放约299.5万，并非本分集播放量；已核对目录与计数，未逐帧观看。\href{https://www.bilibili.com/video/BV1Wv411h7kN/?p=23}{观看分集}。
\end{itemize}
\endgroup
\section{本章小结}
统计学习提供的是关于数据、模型复杂度和未见样本风险的推理框架。它不能保证某个模型在所有工业分布上有效，却能迫使实验者明确假设空间、验证规则、泄漏风险、不确定性和决策成本。

统一界考虑看过数据以后再选模型的影响，嵌套验证把挑选模型和报告成绩分开，配对比较与设备级区间估计则显示结果会怎样波动。比较平均误差前，要先核对任务、抽样单位和尝试方案的预算是否一致；选择告警阈值时，还要看校准和成本。下一章回到原始记录：如果不知道记录何时可用、标签何时截止、样本来自哪台设备，再严谨的统计计算也无法回答原先的维护问题。
